\chapter{2022年北京卷}

\section{单选题}

\begin{problem}
已知全集 \(U=\{x\mid -3<x<3\}\)，集合 \(A=\{x\mid -2<x\le 1\}\)，则 \(C_UA=\)（\quad）.

\choices
    {\((-2,1]\)}
    {\((-3,-2)\cup[1,3)\)}
    {\([-2,1)\)}
    {\((-3,-2]\cup(1,3)\)}
\end{problem}

\begin{answer}
D
\end{answer}



\begin{problem}
若复数 \(z\) 满足 \(\i z=3-4\i\)，则 \(|z|=\)（\quad）.

\choices
    {1}
    {5}
    {7}
    {25}
\end{problem}

\begin{answer}
B
\end{answer}



\begin{problem}
若直线 \(2x+y-1=0\) 是圆 \((x-a)^2+y^2=1\) 的一条对称轴，则 \(a=\)（\quad）.

\choices
    {\(\frac12\)}
    {\(-\frac12\)}
    {1}
    {\(-1\)}
\end{problem}

\begin{answer}
A
\end{answer}



\begin{problem}
已知函数 \(f(x)=\dfrac{1}{1+2^x}\)，则对任意实数 \(x\)，有（\quad）.

\choices
    {\(f(-x)+f(x)=0\)}
    {\(f(-x)-f(x)=0\)}
    {\(f(-x)+f(x)=1\)}
    {\(f(-x)-f(x)=\frac13\)}
\end{problem}

\begin{answer}
C
\end{answer}



\begin{problem}
已知函数 \(f(x)=\cos^2x-\sin^2x\)，则（\quad）.

\choices
    {\(f(x)\) 在 \(\left(-\frac{\pi}{2},-\frac{\pi}{6}\right)\) 上单调递减}
    {\(f(x)\) 在 \(\left(-\frac{\pi}{4},\frac{\pi}{12}\right)\) 上单调递增}
    {\(f(x)\) 在 \(\left(0,\frac{\pi}{3}\right)\) 上单调递减}
    {\(f(x)\) 在 \(\left(\frac{\pi}{4},\frac{7\pi}{12}\right)\) 上单调递增}
\end{problem}

\begin{answer}
C
\end{answer}



\begin{problem}
设 \(\{a_n\}\) 是公差不为 0 的无穷等差数列，则“\(\{a_n\}\) 为递增数列”是“存在正整数 \(N_0\)，当 \(n>N_0\) 时，\(a_n>0\)”的（\quad）.

\choices
    {充分而不必要条件}
    {必要而不充分条件}
    {充分必要条件}
    {既不充分也不必要条件}
\end{problem}

\begin{answer}
C
\end{answer}



\begin{problem}
在北京冬奥会上，国家速滑馆“冰丝带”使用高效环保的二氧化碳跨临界直冷制冰技术，为实现绿色冬奥作出了贡献. 如图描述了一定条件下二氧化碳所处的状态与 \(T\) 和 \(\lg P\) 的关系，其中 \(T\) 表示温度，单位是 \(\mathrm{K}\)；\(P\) 表示压强，单位是 \(\mathrm{bar}\). 下列结论中正确的是（\quad）.

\begin{center}
\bitmapfigure[width=0.34\textwidth]{img/2022/beijing/q7_fig1.png}
\end{center}

\choices
    {当 \(T=220\)，\(P=1026\) 时，二氧化碳处于液态}
    {当 \(T=270\)，\(P=128\) 时，二氧化碳处于气态}
    {当 \(T=300\)，\(P=9987\) 时，二氧化碳处于超临界状态}
    {当 \(T=360\)，\(P=729\) 时，二氧化碳处于超临界状态}
\end{problem}

\begin{answer}
D
\end{answer}



\begin{problem}
若 \((2x-1)^4=a_4x^4+a_3x^3+a_2x^2+a_1x+a_0\)，则 \(a_0+a_2+a_4=\)（\quad）.

\choices
    {40}
    {41}
    {\(-40\)}
    {\(-41\)}
\end{problem}

\begin{answer}
B
\end{answer}



\begin{problem}
已知正三棱锥 \(P-ABC\) 的六条棱长均为 6，\(S\) 是 \(\triangle ABC\) 及其内部的点构成的集合. 设集合 \(T=\{Q\in S\mid PQ\le 5\}\)，则 \(T\) 表示的区域的面积为（\quad）.

\choices
    {\(\frac{3\pi}{4}\)}
    {\(\pi\)}
    {\(2\pi\)}
    {\(3\pi\)}
\end{problem}

\begin{answer}
B
\end{answer}



\begin{problem}
在 \(\triangle ABC\) 中，\(AC=3\)，\(BC=4\)，\(\angle C=90^\circ\). \(P\) 为 \(\triangle ABC\) 所在平面内的动点，且 \(PC=1\)，则 \(\overrightarrow{PA}\cdot \overrightarrow{PB}\) 的取值范围是（\quad）.

\choices
    {\([-5,3]\)}
    {\([-3,5]\)}
    {\([-6,4]\)}
    {\([-4,6]\)}
\end{problem}

\begin{answer}
D
\end{answer}


\section{填空题}

\begin{problem}
函数 \(f(x)=\dfrac{1}{x}+\sqrt{1-x}\) 的定义域是\fillinblank{}.
\end{problem}

\begin{answer}
\((-\infty,0)\cup(0,1]\)
\end{answer}



\begin{problem}
已知双曲线 \(y^2+\dfrac{x^2}{m}=1\) 的渐近线方程为 \(y=\pm \dfrac{\sqrt3}{3}x\)，则 \(m=\)\fillinblank{}.
\end{problem}

\begin{answer}
\(-3\)
\end{answer}



\begin{problem}
若函数 \(f(x)=A\sin x-\sqrt3\cos x\) 的一个零点为 \(\dfrac{\pi}{3}\)，则 \(A=\)\fillinblank{}；\(f\left(\dfrac{\pi}{12}\right)=\)\fillinblank{}.
\end{problem}

\begin{answer}
\(1\)；\(-\sqrt{2}\)
\end{answer}



\begin{problem}
设函数 \(f(x)=\begin{cases}
-ax+1,&x<a,\\
\left(x-2\right)^2,&x\ge a
\end{cases}\)，若 \(f(x)\) 存在最小值，则 \(a\) 的一个取值为\fillinblank{}；\(a\) 的最大值为\fillinblank{}.
\end{problem}

\begin{answer}
\(0\)（答案不唯一）；\(1\)
\end{answer}



\begin{problem}
已知数列 \(\{a_n\}\) 各项均为正数，其前 \(n\) 项和 \(S_n\) 满足 \(a_nS_n=9\)（\(n=1\)，\(2\)，\(\dots\)）. 给出下列四个结论：

\begin{circlelist}
\item \(\{a_n\}\) 的第 2 项小于 3；

\item \(\{a_n\}\) 为等比数列；

\item \(\{a_n\}\) 为递减数列；

\item \(\{a_n\}\) 中存在小于 \(\dfrac{1}{100}\) 的项.
\end{circlelist}

其中所有正确结论的序号是\fillinblank{}.
\end{problem}

\begin{answer}
\(\circled{1}\circled{3}\circled{4}\)
\end{answer}


\section{解答题}

\begin{problem}
在 \(\triangle ABC\) 中，\(\sin 2C=\sqrt3\sin C\).

\begin{enumerate}
\item 求 \(\angle C\)；

\item 若 \(b=6\)，且 \(\triangle ABC\) 的面积为 \(6\sqrt3\)，求 \(\triangle ABC\) 的周长.
\end{enumerate}
\end{problem}

\begin{solution}
\begin{enumerate}
\item 因为 \(C\) 是三角形的内角，所以 \(0<C<\pi\)，从而 \(\sin C>0\). 由
\[
\sin 2C=2\sin C\cos C=\sqrt3\sin C
\]
得 \(2\cos C=\sqrt3\)，所以
\[
\angle C=\dfrac{\pi}{6}.
\]

\item 设 \(a=BC\)，\(c=AB\). 由三角形面积公式，
\[
6\sqrt3=\dfrac12ab\sin C=\dfrac12\cdot a\cdot6\cdot\dfrac12,
\]
所以 \(a=4\sqrt3\). 再由余弦定理，
\[
c^2=a^2+b^2-2ab\cos C
=(4\sqrt3)^2+6^2-2\cdot4\sqrt3\cdot6\cdot\dfrac{\sqrt3}{2}=12,
\]
故 \(c=2\sqrt3\). 因此 \(\triangle ABC\) 的周长为
\[
a+b+c=4\sqrt3+6+2\sqrt3=6+6\sqrt3.
\]
\end{enumerate}
\end{solution}


\begin{problem}
如图，在三棱柱 \(ABC-A_1B_1C_1\) 中，侧面 \(BCC_1B_1\) 为正方形，平面 \(BCC_1B_1\perp\) 平面 \(ABB_1A_1\)，\(AB=BC=2\)，\(M\)，\(N\) 分别为 \(A_1B_1\)，\(AC\) 的中点.



\begin{enumerate}
\item 求证：\(MN\parallel\) 平面 \(BCC_1B_1\)；

\item 再从下列两个条件中选择一个作为已知，求直线 \(AB\) 与平面 \(BMN\) 所成角的正弦值.

\begin{circlelist}
\item \(AB\perp MN\)；

\item \(BM=MN\).
\end{circlelist}

\end{enumerate}

\bitmapfigure[width=0.30\textwidth]{img/2022/beijing/q17_fig1.png}
\end{problem}


\begin{solution}
\begin{enumerate}
\item 以 \(B\) 为原点，分别以 \(BC\) 和 \(BB_1\) 所在直线为 \(x\) 轴和 \(z\) 轴，建立空间直角坐标系. 由侧面 \(BCC_1B_1\) 为边长 \(2\) 的正方形，且两个平面垂直，可取
\[
B(0,0,0),\quad C(2,0,0),\quad B_1(0,0,2),\quad C_1(2,0,2).
\]
平面 \(ABB_1A_1\) 为平面 \(x=0\)，设 \(A=(0,u,v)\). 由 \(AB=2\) 得 \(u^2+v^2=4\)，又因为是三棱柱，所以 \(A_1=(0,u,v+2)\). 因而
\[
M=\left(0,\dfrac{u}{2},\dfrac{v}{2}+2\right),\qquad N=\left(1,\dfrac{u}{2},\dfrac{v}{2}\right).
\]
所以 \(\overrightarrow{MN}=(1,0,-2)\). 而平面 \(BCC_1B_1\) 的方程为 \(y=0\)，故 \(MN\parallel\) 平面 \(BCC_1B_1\).

\item 令 \(h=\dfrac{u}{2}\)，\(w=\dfrac{v}{2}\)，则 \(h^2+w^2=1\)，且
\[
M=(0,h,w+2),\qquad N=(1,h,w).
\]
取平面 \(BMN\) 的一个法向量
\[
\bs n=\overrightarrow{BM}\times\overrightarrow{BN}=(-2h,w+2,-h).
\]
直线 \(AB\) 的方向向量可取 \(\overrightarrow{BA}=(0,2h,2w)\)，于是直线 \(AB\) 与平面 \(BMN\) 所成角 \(\theta\) 满足
\[
\sin\theta=\dfrac{\left|\overrightarrow{BA}\cdot\bs n\right|}{\left|\overrightarrow{BA}\right|\left|\bs n\right|}
=\dfrac{2|h|}{\sqrt{5h^2+(w+2)^2}}.
\]

若选择条件\(\circled{1}\)，则
\[
\overrightarrow{BA}\cdot\overrightarrow{MN}=(0,2h,2w)\cdot(1,0,-2)=-4w=0,
\]
所以 \(w=0\)，再由 \(h^2+w^2=1\) 得 \(h^2=1\). 因此
\[
\sin\theta=\dfrac{2}{\sqrt{5+4}}=\dfrac{2}{3}.
\]

若选择条件\(\circled{2}\)，则
\[
|BM|^2=h^2+(w+2)^2=5+4w,
\qquad |MN|^2=1+4=5.
\]
由 \(BM=MN\) 得 \(w=0\)，再由 \(h^2+w^2=1\) 得 \(h^2=1\). 同理
\[
\sin\theta=\dfrac{2}{3}.
\]
因此，在两种条件下所求正弦值均为 \(\dfrac{2}{3}\).
\end{enumerate}
\end{solution}


\begin{problem}
在校运动会上，只有甲，乙，丙三名同学参加铅球比赛，比赛成绩达到 \(9.50\text{ m}\) 以上（含 \(9.50\text{ m}\)）的同学将获得优秀奖. 为预测获得优秀奖的人数及冠军得主，收集了甲，乙，丙以往的比赛成绩，并整理得到如下数据（单位：\(\mathrm{m}\)）：

甲：\(9.80\)，\(9.70\)，\(9.55\)，\(9.54\)，\(9.48\)，\(9.42\)，\(9.40\)，\(9.35\)，\(9.30\)，\(9.25\)；

乙：\(9.78\)，\(9.56\)，\(9.51\)，\(9.36\)，\(9.32\)，\(9.23\)；

丙：\(9.85\)，\(9.65\)，\(9.20\)，\(9.16\).

假设用频率估计概率，且甲，乙，丙的比赛成绩相互独立.

\begin{enumerate}
\item 估计甲在校运动会铅球比赛中获得优秀奖的概率；

\item 设 \(X\) 是甲，乙，丙在校运动会铅球比赛中获得优秀奖的总人数，估计 \(X\) 的数学期望 \(E(X)\)；

\item 在校运动会铅球比赛中，甲，乙，丙谁获得冠军的概率估计值最大？（结论不要求证明）.
\end{enumerate}
\end{problem}


\begin{solution}
\begin{enumerate}
\item 甲以往的 \(10\) 次成绩中，有 \(4\) 次达到 \(9.50\text{ m}\) 以上，因此估计甲获得优秀奖的概率为
\[
P_A=\dfrac{4}{10}=0.4.
\]

\item 乙和丙获得优秀奖的概率估计值分别为
\[
P_B=\dfrac{3}{6}=\dfrac12,\qquad P_C=\dfrac{2}{4}=\dfrac12.
\]
令 \(I_A\)，\(I_B\)，\(I_C\) 分别表示甲，乙，丙获得优秀奖的指示变量，则
\[
X=I_A+I_B+I_C.
\]
由数学期望的线性性质，
\[
E(X)=E(I_A)+E(I_B)+E(I_C)=P_A+P_B+P_C
=0.4+\dfrac12+\dfrac12=\dfrac75.
\]

\item 三人的历史最好成绩分别为 \(9.80\text{ m}\)，\(9.78\text{ m}\)，\(9.85\text{ m}\)，其中丙的历史最好成绩最高，因此估计丙获得冠军的概率最大.
\end{enumerate}
\end{solution}


\begin{problem}
已知椭圆 \(E:\frac{x^2}{a^2}+\frac{y^2}{b^2}=1\)（\(a>b>0\)）的一个顶点为 \(A(0,1)\)，焦距为 \(2\sqrt3\).

\begin{enumerate}
\item 求椭圆 \(E\) 的方程；

\item 过点 \(P(-2,1)\) 作斜率为 \(k\) 的直线与椭圆 \(E\) 交于不同的两点 \(B\)，\(C\)，直线 \(AB\)，\(AC\) 分别与 \(x\) 轴交于点 \(M\)，\(N\). 当 \(|MN|=2\) 时，求 \(k\) 的值.
\end{enumerate}
\end{problem}


\begin{solution}
\begin{enumerate}
\item 椭圆的短半轴为 \(b\)，由顶点 \(A(0,1)\) 得 \(b=1\). 设半焦距为 \(c\)，由焦距为 \(2\sqrt3\) 得 \(c=\sqrt3\). 又 \(a^2=b^2+c^2\)，所以 \(a^2=4\). 因而椭圆 \(E\) 的方程为
\[
\dfrac{x^2}{4}+y^2=1.
\]

\item 设过点 \(P(-2,1)\) 的直线为
\[
y-1=k(x+2),
\]
即 \(y=kx+2k+1\). 将其代入椭圆方程，得交点横坐标 \(x_1\)，\(x_2\) 满足
\[
(1+4k^2)x^2+8k(2k+1)x+16k(k+1)=0.
\]
因为交于不同的两点，判别式为
\[
\Delta=\left[8k(2k+1)\right]^2-4(1+4k^2)\cdot16k(k+1)=-64k>0,
\]
所以 \(k<0\). 由韦达定理，
\[
x_1+x_2=-\dfrac{8k(2k+1)}{1+4k^2},\qquad x_1x_2=\dfrac{16k(k+1)}{1+4k^2}.
\]
从而
\[
(x_1+2)(x_2+2)=x_1x_2+2(x_1+x_2)+4=\dfrac{4}{1+4k^2}.
\]

设相应的交点为 \(B(x_1,y_1)\)，\(C(x_2,y_2)\). 过 \(A(0,1)\) 与 \((x_i,y_i)\) 的直线在 \(x\) 轴上的截距记为 \(u_i\)，则
\[
u_i=\dfrac{x_i}{1-y_i}=-\dfrac{x_i}{k(x_i+2)}.
\]
因此
\[
|MN|=|u_1-u_2|=\dfrac{2|x_1-x_2|}{|k|(x_1+2)(x_2+2)}.
\]
又由判别式得 \(|x_1-x_2|=\dfrac{\sqrt{\Delta}}{1+4k^2}=\dfrac{8\sqrt{-k}}{1+4k^2}\). 代入上式得
\[
|MN|=\dfrac{4\sqrt{-k}}{|k|}=\dfrac{4}{\sqrt{-k}}.
\]
由 \(|MN|=2\)，得 \(\sqrt{-k}=2\)，所以 \(k=-4\).
\end{enumerate}
\end{solution}


\begin{problem}
已知函数 \(f(x)=\e^x\ln(1+x)\).

\begin{enumerate}
\item 求曲线 \(y=f(x)\) 在点 \((0,f(0))\) 处的切线方程；

\item 设 \(g(x)=f'(x)\)，讨论函数 \(g(x)\) 在 \([0,+\infty)\) 上的单调性；

\item 证明：对任意的 \(s,t\in(0,+\infty)\)，有 \(f(s+t)>f(s)+f(t)\).
\end{enumerate}
\end{problem}


\begin{solution}
\begin{enumerate}
\item 由
\[
f(0)=0,\qquad f'(x)=\e^x\left(\ln(1+x)+\dfrac{1}{1+x}\right),
\]
得 \(f'(0)=1\). 因此曲线在点 \((0,f(0))=(0,0)\) 处的切线方程为
\[
y-0=1(x-0),
\]
即 \(y=x\).

\item 由 \(g(x)=f'(x)\)，得
\[
g'(x)=f''(x)=\e^x\left(\ln(1+x)+\dfrac{2}{1+x}-\dfrac{1}{(1+x)^2}\right).
\]
令 \(H(x)=\ln(1+x)+\dfrac{2}{1+x}-\dfrac{1}{(1+x)^2}\). 当 \(x\geq0\) 时，
\[
H(0)=0,\qquad H'(x)=\dfrac{x^2+1}{(1+x)^3}>0.
\]
所以 \(H(x)\geq0\)，从而 \(g'(x)\geq0\). 因此 \(g(x)\) 在 \([0,+\infty)\) 上单调递增.

\item 固定 \(s\in(0,+\infty)\)，令
\[
F(t)=f(s+t)-f(s)-f(t),\qquad t\geq0.
\]
因为 \(f(0)=0\)，所以 \(F(0)=0\). 又由第（2）问知 \(g\) 在 \([0,+\infty)\) 上严格递增，故对 \(t\geq0\)，
\[
F'(t)=g(s+t)-g(t)>0.
\]
所以当 \(t>0\) 时，\(F(t)>F(0)=0\)，即
\[
f(s+t)>f(s)+f(t).
\]
\end{enumerate}
\end{solution}


\begin{problem}
已知 \(Q:a_1\)，\(a_2\)，\(\dots\)，\(a_k\) 为有穷整数数列. 给定正整数 \(m\)，若对任意的 \(n\in\{1,2,\dots,m\}\)，在 \(Q\) 中存在 \(a_i\)，\(a_{i+1}\)，\(\dots\)，\(a_{i+j}\)（\(j\ge 0\)），使得 \(a_i+a_{i+1}+\dots+a_{i+j}=n\)，则称 \(Q\) 为 \(m\)-连续可表数列.

\begin{enumerate}
\item 判断 \(Q:2\)，\(1\)，\(4\) 是否为 5-连续可表数列？是否为 6-连续可表数列？说明理由；

\item 若 \(Q:a_1\)，\(a_2\)，\(\dots\)，\(a_k\) 为 8-连续可表数列，求证：\(k\) 的最小值为 4；

\item 若 \(Q:a_1\)，\(a_2\)，\(\dots\)，\(a_k\) 为 20-连续可表数列，且 \(a_1+a_2+\dots+a_k<20\)，求证：\(k\ge 7\).
\end{enumerate}
\end{problem}

\begin{solution}
\begin{enumerate}
\item 由定义，数列 \(Q:2,1,4\) 的非空连续子段的和为
\[
2,\quad1,\quad4,\quad2+1=3,\quad1+4=5,\quad2+1+4=7.
\]
其中包含 \(1,2,3,4,5\)，所以 \(Q:2,1,4\) 是 \(5\)-连续可表数列；但这些连续子段的和中没有 \(6\)，所以它不是 \(6\)-连续可表数列.

\item 长度为 \(k\) 的数列至多有
\[
1+2+\cdots+k=\dfrac{k(k+1)}{2}
\]
个非空连续子段. 若 \(k\leq3\)，则连续子段数不超过 \(6\)，不可能表示 \(1,2,\dots,8\) 这 \(8\) 个互不相同的整数，所以 \(k\geq4\).

取数列 \(Q:1,1,3,3\)，其连续子段的和中有
\[
1,\quad1+1=2,\quad3=3,\quad1+1+3=5,\quad1+3=4,\quad3+3=6,\quad1+3+3=7,\quad1+1+3+3=8.
\]
再由上述连续子段的和可知 \(1,2,\dots,8\) 均可表示. 因而 \(k=4\) 时确实存在 \(8\)-连续可表数列，\(k\) 的最小值为 \(4\).

\item 若 \(k\leq5\)，非空连续子段数至多为 \(15\)，小于需要表示的 \(20\) 个整数，所以只需排除 \(k=6\).

假设 \(k=6\). 此时共有 \(\dfrac{6\cdot7}{2}=21\) 个非空连续子段，其中整个数列作为一个连续子段，其和满足
\[
a_1+a_2+\cdots+a_6<20,
\]
不能表示整数 \(20\). 因此其余 \(20\) 个真子段必须恰好表示 \(1,2,\dots,20\)，每个真子段的和都是正整数.

特别地，每个单项子段都是上述真子段，所以 \(a_i>0\)（\(i=1,2,\dots,6\)）. 于是任意连续子段的和都不超过全体各项的和，而全体各项的和小于 \(20\)，这与必须表示 \(20\) 矛盾. 所以 \(k\ne6\)，从而 \(k\geq7\).
\end{enumerate}
\end{solution}
